\chapter{Blocks 11--12: AI, Right Questions \& Capstone}

\textbf{Primary:} AI prompting, data expectations, reproducible pipelines, capstone project

\begin{brainbox}[The meta-skill]
Tools change. Methods change. The ability to ask the right question, recognize unexpected data, and evaluate AI output --- that lasts your whole career.
\end{brainbox}


\section{Block 11: Asking the Right Questions + AI Mastery}

\sprint{1}{The Question Framework}{25}
\begin{itemize}[nosep]
  \item Before EVERY analysis, write down answers to: What's the biological question? What would a yes look like? What would a no look like?
  \item Data inventory: how many rows? columns? What does each represent? What are the expected ranges?
  \item Sanity checks: $\sim$4--5M variants per human genome, $\sim$20K protein-coding genes, scRNA-seq: 500--5000 genes/cell
  \item Red flags to watch for: flat Manhattan plot (no signal), bimodal quality scores, batch-correlated PCA
\end{itemize}
\begin{winbox}[Checkpoint]
You wrote a pre-analysis checklist for a dataset. This prevents weeks of wasted analysis on broken data.
\end{winbox}

\subsection{The Question Map}

\begin{center}
\begin{tabularx}{\textwidth}{l X}
\toprule
\headertext{Stage} & \headertext{Ask This} \\
\midrule
Before you start & What is the question? What would a definitive answer look like? \\
Loading data & What do rows/columns represent? Expected dimensions? Missing values? \\
Choosing methods & What does this method assume? Does my data meet the assumptions? \\
After QC & Do distributions look right? Do controls work? Any batch effects? \\
Interpreting results & Is the effect size meaningful? Could confounders explain this? \\
Before publishing & Can someone else reproduce this? Are all parameters recorded? \\
\bottomrule
\end{tabularx}
\end{center}

\sprint{2}{Prompting AI Like a Bioinformatician}{25}
\begin{itemize}[nosep]
  \item Context-first prompting: ``I have a 10x scRNA-seq h5ad file, 50K cells, human PBMCs. I want to...''
  \item Constraint prompting: ``Use scanpy 1.9+, avoid deprecated functions, include comments at each QC threshold''
  \item Verification prompting: ``Include assertions that check the output dimensions and warn if something unexpected happens''
  \item Debugging: paste the FULL error message + stack trace. Never describe the error in your own words.
\end{itemize}
\begin{winbox}[Checkpoint]
You wrote a prompt that generated working, well-documented bioinformatics code on the first try.
\end{winbox}

\sprint{3}{Spotting AI Hallucinations}{25}
\begin{itemize}[nosep]
  \item Common AI mistakes: inventing function names, confusing package versions, incorrect parameter names
  \item Always test: run AI code on a small known-answer dataset before trusting it on real data
  \item Cross-validate: ask two different AI tools the same question. If they disagree, investigate.
  \item Version trap: AI may mix up scanpy 1.8 vs 1.10 syntax, or old Seurat v3 vs v5 functions
\end{itemize}
\begin{winbox}[Checkpoint]
You found and corrected an AI hallucination. This is a critical skill for AI-augmented research.
\end{winbox}

\begin{warnbox}[AI ground rules for research]
\begin{itemize}[nosep]
  \item Always run AI-generated code before trusting it.
  \item Cite AI assistance in your methods section --- transparency matters.
  \item Never use AI for novel method development without deep verification.
  \item AI is a study partner, not a replacement for understanding.
\end{itemize}
\end{warnbox}


\section{Block 12: Reproducible Pipeline + Capstone}

\sprint{4}{Snakemake: Your Pipeline in Code}{50}
\begin{itemize}[nosep]
  \item Install: \texttt{mamba install -c bioconda -c conda-forge snakemake}. Docs: \url{https://snakemake.readthedocs.io/}.
  \item Snakemake defines a pipeline as a set of rules. Each rule: input files $\to$ output files + a shell/Python command.
  \item Start simple: rule download + rule qc + rule align. Run with \texttt{snakemake --cores 4}.
  \item Snakemake automatically figures out the execution order. Change an input file? It re-runs only what's needed.
  \item Add a \texttt{conda:} directive to each rule for per-step environments. Now it's fully reproducible.
  \item Alternative: Nextflow + nf-core if you prefer. Install Nextflow with \texttt{curl -s https://get.nextflow.io | bash}. nf-co.re (\url{https://nf-co.re}) has community-curated pipelines for most assays.
\end{itemize}
\begin{winbox}[Checkpoint]
You have a multi-step Snakemake pipeline that someone else could clone and run.
\end{winbox}

\subsection{Capstone Project}

This is your final project. It ties together everything from the past 11 blocks into a single, complete analysis.

\begin{dobox}[What to build]
\begin{itemize}[nosep]
  \item Choose a biological question and a public dataset (GEO, SRA, GWAS Catalog).
  \item Build a reproducible pipeline (Snakemake or Nextflow) in a Git repository.
  \item Analyze with Python and/or R, using appropriate statistical methods.
  \item Produce at least 3 publication-quality figures with captions.
  \item Write a 5--10 page report: question, methods, results, interpretation.
  \item Include an AI journal: how you used AI tools, what worked, what didn't.
\end{itemize}
\end{dobox}

\begin{tipbox}[Suggested capstone topics]
\begin{itemize}[nosep]
  \item Complete Problem Set 4 (QTL Mapping) from the Running Project --- this is the recommended capstone if you've been working through Problem Sets 1--3
  \item Reproduce a published scRNA-seq analysis with updated methods (new normalization, better clustering)
  \item GWAS-to-function: fine-map a locus, colocalize with eQTLs, identify candidate genes
  \item Cross-ancestry PRS: build polygenic risk scores, test portability across populations
  \item CRISPR screen reanalysis: apply MAGeCK + BAGEL2 to a published dataset, compare hit lists
\end{itemize}
\end{tipbox}

\sprint{5}{The Reproducibility Test}{25}
\begin{itemize}[nosep]
  \item Clone your repo to a different computer (or ask a friend/AI to try).
  \item Can they install the environment and run the pipeline end-to-end?
  \item If yes: you're done. If no: fix what's broken. This is the real test.
\end{itemize}
\begin{winbox}[Checkpoint]
Someone else can reproduce your analysis from your repo. That is the standard to aim for.
\end{winbox}

\begin{restbox}[End of the curriculum]
At this point you can process genomic data in Unix, analyze it in Python and R, apply rigorous statistics, and produce reproducible research. The foundation is in place --- continued learning comes from applying these skills to real research questions.
\end{restbox}
